\addtocontents{toc}{\protect\setcounter{tocdepth}{0}}
\chapter[Conclusiones]{Conclusiones}

\section{Contrastación de objetivos}

El objetivo general del proyecto es desarrollar una solución móvil que
centralice el registro, el seguimiento, la consulta y la orientación del
cuidador. A la fecha de este informe, OnCuidar cumple ese propósito en su flujo
principal. El cuidador crea su cuenta, registra la valoración diaria del
paciente, recibe un nivel de alerta, consulta y exporta el historial, resuelve
dudas en la biblioteca, el asistente o las preguntas frecuentes y programa
recordatorios.

Respecto de los objetivos específicos, el análisis de las necesidades del
cuidador se concretó en 29 requerimientos funcionales derivados de las historias
de usuario y de los contenidos de Enfermería. El diseño se materializó en una
arquitectura por capas en la aplicación, una API REST propia con un modelo de
datos relacional y cifrado en el cliente, y una interfaz validada con
prototipos. La implementación completó los módulos de
los sprints~1 y~2, y el sprint~3 aborda la separación de datos por paciente, el
aislamiento entre usuarios y el funcionamiento sin conexión. La validación
técnica se apoya en 294 pruebas automatizadas de la aplicación y 54 del
backend, que se ampliarán con los requerimientos del último sprint.

Los objetivos del software, que son ordenar el seguimiento domiciliario y
orientar al cuidador ante signos de alerta, se alcanzarán con el uso de la
aplicación por parte de las familias, más allá del término del proyecto.

\section{Conclusiones sobre herramientas y metodología}

Flutter resultó adecuado para un equipo de dos personas con un plazo acotado,
pues una sola base de código cubre Android e iOS. Por requerimiento de la
universidad, el backend se construyó como una API REST propia con Node.js y una
base de datos relacional PostgreSQL, empaquetadas con Docker Compose y destinadas
a ejecutarse como contenedores Podman en el servidor institucional. Esta decisión
entrega control total sobre los datos y su cifrado y cumple lo solicitado, a
cambio de asumir la operación del servidor, la gestión de las claves, el envío
de correos y el mantenimiento de la base de datos, tareas que un servicio
administrado resuelve por el equipo. Dos decisiones técnicas destacan por su
valor en un contexto de salud. El cifrado de los datos sensibles en el propio
dispositivo permite que ni siquiera el servidor ni la base de datos los tengan
legibles, y el servidor rechaza los campos clínicos en claro. La orientación
basada en reglas y respuestas validadas, en lugar de inteligencia artificial
generativa, mantiene cada respuesta dentro del alcance definido por Enfermería.

Scrum permitió mostrar una versión funcional a la contraparte al final de cada
sprint y ajustar textos y criterios a tiempo. La estimación por puntos de caso
de uso mostró que el alcance completo excedía el presupuesto de horas, lo que
obligó a priorizar el backlog desde el inicio.

Como trabajo futuro, la primera tarea es completar el despliegue en el servidor
de la universidad una vez activada la cuenta, y servir la API sobre HTTPS. Luego
corresponde programar respaldos periódicos de PostgreSQL, validar la aplicación
en dispositivos reales y en iOS, que aún no se ha probado, y ejecutar el plan de
certificación del funcionamiento sin conexión.

\section{Conclusiones generales}

Desde el punto de vista académico, el proyecto integra ingeniería de
requerimientos, diseño de software, bases de datos, seguridad de la información
y pruebas en un problema real del sistema de salud. Desde el punto de vista
personal, trabajar con una contraparte clínica enseñó que en salud cada texto y
cada umbral importan, y que un buen software de apoyo reconoce sus límites y
deriva al profesional cuando corresponde.
